\section{Method}
% Setup, base BA objective, then the superquadric surface prior (the contribution), association/EM, and prior registration.
We are given a sparse set of $N$ views of an indoor scene with known intrinsics and want accurate camera poses. The pipeline runs in two stages: a feed-forward stage that produces an initial reconstruction, and a bundle adjustment stage that refines the camera poses, with the superquadric prior entering as an extra cost in the second stage.

\paragraph{Initialization and correspondences}
We initialize the cameras and scene geometry with VGGT \cite{wang2025vggtvisualgeometrygrounded} through the MapAnything harness \cite{keetha2026mapanythinguniversalfeedforwardmetric}, which gives a camera pose and a dense depth map per view in a single forward pass. To get the 2D--2D matches that bundle adjustment optimizes against, we run MASt3R \cite{leroy2024groundingimagematching3d} on every ordered view pair and keep its dense reciprocal correspondences above a confidence threshold. Each surviving match is unprojected to a camera-frame ray with the VGGT intrinsics, rotated into the world frame using the VGGT poses, and triangulated by closest-approach midpoint of the two rays. We discard near-parallel rays, and we keep a triangulated point only if it has positive depth in both cameras and reprojects inside both images. This yields the initial 3D points $X_j$ and their observations $x_{ij}$.

\paragraph{Bundle adjustment objective}
Each camera $i$ is parameterized by a world-to-camera rotation $R_i$ and translation $t_i$; the intrinsics are held fixed during optimization to avoid the focal-length/depth ambiguity. The base cost is the robustified reprojection error over all observations,
\begin{equation}
E_{\text{reproj}} = \sum_{i,j} \rho\!\left( \lVert \pi(R_i, t_i, X_j) - x_{ij} \rVert^2 \right),
\end{equation}
where $\pi$ is the pinhole projection and $\rho$ is the Huber loss with threshold $\delta$ (default $2.0$ px, set to $1.0$ px in our reported runs), which down-weights gross MASt3R mismatches. We solve this least-squares problem with Ceres using a sparse Schur factorization, optimizing cameras and points jointly. The gauge is fixed by holding the first camera constant, which removes the $6$-DoF freedom of the global frame.

\paragraph{Superquadric surface prior}
Reprojection alone under-constrains the cameras when views barely overlap. We add a structural prior that pulls the reconstructed points onto the known scene geometry, expressed as a set of superquadrics \cite{Paschalidou2019CVPR, fedele2026superdec3dscenedecomposition} fit offline (see below). Each point $X_j$ is associated to one superquadric $s(j)$ and mapped into that primitive's canonical frame, $q_j = R_{s(j)}^{\top}(X_j - t_{s(j)})$, where $R_s$ and $t_s$ are the primitive's rotation and centre. The inside--outside function of a superquadric with half-extents $a_1, a_2, a_3$ and shape exponents $\varepsilon_1, \varepsilon_2$ is
\begin{equation}
F(q) = \left( \left|\tfrac{q_x}{a_1}\right|^{2/\varepsilon_2} + \left|\tfrac{q_y}{a_2}\right|^{2/\varepsilon_2} \right)^{\varepsilon_2/\varepsilon_1} + \left|\tfrac{q_z}{a_3}\right|^{2/\varepsilon_1},
\end{equation}
in the inverse Solina form \cite{article}, with $F^{-\varepsilon_1/2}-1$ positive inside the surface, negative outside, and zero exactly on it. The full objective adds a weighted one-sided surface term to the reprojection cost:
\begin{equation}
E = E_{\text{reproj}} + \sum_{j} \Big( \lambda\, \lVert q_j \rVert\, \max\!\big(0,\, 1 - F(q_j)^{-\varepsilon_1/2} \big) \Big)^2 .
\end{equation}
The $\max(0,\cdot)$ makes the term a hinge: it penalizes only points that lie \emph{outside} their associated surface and leaves interior points free, so the prior pulls points that have drifted off the scene back onto it and leaves points already inside a primitive alone. The factor $\lVert q_j \rVert$ turns the unitless inside--outside value into an approximate radial distance in meters. The weight $\lambda$ is a pixels-per-meter factor that puts the surface residual on the same scale as the reprojection residual; setting $\lambda=0$ recovers plain reprojection BA. We use $\lambda=15$ and a separate Huber of $2.749$ on the surface term. Both costs are summed into one Ceres problem and minimized together.

\paragraph{Association and EM refinement}
Each point is bound to its nearest superquadric by radial distance, $r(p) = \lVert q \rVert \, \lvert F(q)^{-\varepsilon_1/2} - 1 \rvert$, taking the argmin over the active primitives. A point whose nearest-primitive distance exceeds a cutoff (we use $0.0372$ m) is left unassigned and contributes no surface term, which keeps stray triangulations from being pulled toward unrelated geometry. Because Ceres moves the points during optimization, a fixed association computed at the start drifts out of date. We therefore re-estimate it in an EM-style loop: a reprojection-only warm-up solve first cleans up the structure so the first association is reliable, then we alternate an E-step that re-associates the current points with an M-step that runs a short surface BA solve. We use two outer iterations of $41$ inner Ceres iterations each. With one outer iteration this reduces to associate-once, solve-once.

\paragraph{Prior registration}
The superquadrics are fit offline by SuperDec \cite{fedele2026superdec3dscenedecomposition} on the clean, full-scene point cloud back-projected from ground-truth depth and instance masks, and they live in the scene's world frame. Our setting is camera-pose refinement in a \emph{known} environment, so this geometry is available as a fixed prior rather than test-time information about the cameras we solve for. To use it, we bring the superquadrics from the scene frame into the reconstruction frame produced by VGGT. We fit a $\mathrm{Sim}(3)$ from the ground-truth camera centres to the predicted camera centres with Umeyama alignment, invert it, and apply the inverse to the primitives before association. This step uses ground-truth camera centres only to register the prior; the camera \emph{poses} are still what bundle adjustment refines.
